\subsection{正三次元対角chartの共鳴は$a^3=-1$に限る}
\label{subsec:h3-positive3-diagonal-top}
\begin{theorem}\label{thm:h3-positive3-diagonal-top}
$a\ne0$に対し、次の先頭環を考える：
\[
 B_0=k\langle x,y,v,z\rangle/
 (xy-a^2yx,\ zx-axz,\ zy-a^{-1}yz,
 vx-axv,\ vy+a^{-1}yv,\ v^2-z^2).
\]
これはnoetherian quadratic Koszul AS4型14641の整域で、
$\mu_{B_0}=(x,y,v,-z)$である。正規三次の全核は
\[
 W_1(B_0)=
 \begin{cases}
 0,&a^3\ne-1,\\
 ky^3,&a^3=-1,\ a\ne-1,\\
 y\langle x^2,y^2,v^2\rangle,&a=-1.
 \end{cases}
\]
従ってこの先頭族の全三次切断はH3でない。
さらに、実際の中山作用が$(x,y,v,-z)$である14641包絡の表示が
\begin{align*}
 xy&=a^2yx,&zx&=axz,&zy&=a^{-1}yz,\\
 vx&=axv+d_x(W),&vy&=-a^{-1}yv+d_y(W),\\
 v^2-z^2&=\ell(W)v+q_2(W)
\end{align*}
なら、$a^3\ne-1$の全係数で$W_1(B)=0$であり、H3は存在しない。
ここで$d_x,d_y,q_2$は任意の二次式、$\ell$は任意の一次式である。
共鳴値での先頭三次の因子化は、補正後の三次の因子化を意味しない。
\end{theorem}
\begin{proof}
$r=v+z,s=v-z$と置くと$A=k\langle r,s\rangle/(rs+sr)$上の二段Ore表示
\[
 B_0=A[x;\sigma_x][y;\sigma_y],\quad
 \sigma_x(r,s)=a^{-1}(r,s),\quad
 \sigma_y(r,s,x)=(-as,-ar,a^{-2}x)
\]
を得る。従って$r^is^jx^ky^\ell$がPBW基底であり、
noetherian、整域、Koszul AS4である。
Ore中山公式\cite[Theorem 3.3]{LWWOre}を適用する。
$\mu_A=(-r,-s)$、$\operatorname{hdet}_A\sigma_x=a^{-2}$なので
第一拡大の中山は$(-ar,-as,a^{-2}x)$。
第二拡大では$\sigma_y^{-1}$を合成して$(r,s,x)\mapsto(s,r,x)$となる。
$\sigma_y$は$A$の二次関係を$a^2$倍し$x$を$a^{-2}$倍するため、
第一拡大でのhomological determinantは1、従って$y$は固定される。
これで表示した実際の中山作用を得る。

正常三次$G$は中山固定である。
外側平面$A$の次数$n=0,1,2,3$で、中山固定部分の次元は$1,1,1,2$。
それぞれ$1,v,v^2$、および二次元の三次固定空間である。
従って$G$を$W$次数$(i,j)$と外側次数$n$、$i+j+n=3$で分解する。
$Gx=xG,Gy=yG$の必要条件は各部分で
\[
 a^{n-2j}=1,\qquad (-1)^n a^{2i-n}=1.
\]
$n=1$の三場合と$n=3$では、この二式のLaurent多項式の最大公約数は1なので全て零となる。
残りは$x^iy^jv^n$、$n=0,2$の六つの単項式である。
さらに$Gv=vG,Gz=-zG$を課すと
\[
 (-1)^j a^{i-j}=1,\qquad a^{i-j}=-1
\]
が必要十分となる。従って$j$は奇数で、残る候補は
$y^3,x^2y,yv^2$だけ。
最初は$a^3=-1$、後二者は$a=-1$が必要十分である。
異なる$W$次数の成分は独立で、同じ$W$次数内の固定外側基底もPBWで独立なので、
この消去は係数の相殺を残さない。各候補の四正常化式を直接確認すると、全核が得られる。
全て$y$で因子化するので商は整域でない。

補正を持つ表示について、$x,y$の重みを1、$v,z$の重みを2とする。
先頭関係から$B_0\twoheadrightarrow\operatorname{gr}B$があり、
各元の元来の斉次次数を固定した有限filtrationの次元を比較できる。
$B$は14641 Hilbert値を持ち、$B_0$も同じなので、この全射は同型である。
正常三次$G\ne0$の重み先頭項は非零で、同じ四正常化式と中山固定条件を満たす。
従って$\operatorname{in}(G)\in W_1(B_0)$。
$a^3\ne-1$ではこの核が零なので$G$は存在しない。
共鳴値ではこの議論は先頭項を制限するだけで、変形した$G$の因子化を主張しない。
\end{proof}

\texttt{verify\_cubic\_positive3\_diagonal\_top.py}は二段Oreの全生成元三重積、
元の六関係、外側固定部分の次元、十一係数全体の消去に用いるLaurent最大公約数、
共鳴候補の正常化式を厳密検算する。
反対角作用、非半単純作用、$\gamma\ne0$やcross射影退化をこの表示で被覆したとはしない。

\begin{theorem}[正三次元対角normalizing chartの全排除]
\label{thm:h3-positive3-diagonal-full-exclusion}
前定理の補正を持つ表示に属するH3は、共鳴値$a^3=-1$も含めて存在しない。
$a=-1$では$v$を$W$だけ平行移動した後、
\[
 vx=-xv+\kappa y^2,\quad vy=yv,\quad
 v^2-z^2=A x^2+Bxy+C y^2
\]
となり、全正常三次は$y\langle x^2,y^2,v^2\rangle$である。
原始6乗根では全補正が除去され、$W_1=ky^3$となる。
表示した任意の補正係数についてAS4の存在を主張するのではなく、
H3に必要なAS4・整域性・中山作用を満たす表示を全て排除する。
\end{theorem}
\begin{proof}
非共鳴値は前定理によるので$a^3=-1$とする。
$T=(xy=a^2yx)$、$g(x)=ax,g(y)=-a^{-1}y$と置く。
$d_x=d_0x^2+d_1xy+d_2y^2$、$d_y=e_0x^2+e_1xy+e_2y^2$は
$g$-skew derivationであり、その条件は
\begin{gather*}
 a(1-a)e_0=(1+a)d_2=0,\\
 (1+a^{-3})d_0+(a-1)e_1=0,\qquad
 (1+a^{-1})d_1+(a-a^{-2})e_2=0.
\end{gather*}
これは$g(x)d_y+d_xy-a^2(g(y)d_x+d_yx)=0$の四係数である。
$v$の$W$平行移動による内部derivationの二列は
\[
 (d_0,e_1)=(1-a,1+a^{-3}),\qquad
 (d_1,e_2)=(a^{-2}-a,1+a^{-1}).
\]

$a=-1$では$e_0=e_1=e_2=0$で、$d_0,d_1$は内部移動で消える。
従って$d_x=\kappa y^2,d_y=0$となる。
特に$y$は$x,v$と可換、$z$と反可換な正規元であり、
$x^2,y^2,v^2$は$x,y,v$と可換である。
$v^2,z^2$を$x$と比較すると、外側平方関係から
\[
 \ell(W)[v,x]=0
\]
となる。もし$\ell\ne0$なら整域性から$[v,x]=0$、従って
$2xv=\kappa y^2$となり、$v$と$z^2$が$T$の斉次商体に入る。
これは補題\ref{lem:h3-positive3-rational-collapse}に反する。
よって$\ell=0$で、表示した平方関係を得る。
その右辺は$z$の反転作用で固定されるから、$v^2$は$z$とも可換である。
従って$x^2,y^2,v^2$は全て中心に入り、
$yx^2,y^3,yv^2$は正常化作用$\mu_B$を持つ独立な三次である。
前定理の重み先頭環により$\dim W_1(B)\le\dim W_1(B_0)=3$なので、
この三つが全正常核を張る。全て$y$で因子化して非H3となる。

原始6乗根の場合は$a^2-a+1=0$である。
上の四条件の階数は4、内部二列は独立であり、全skew derivationは内部である。
従って$v$の平行移動で$d_x=d_y=0$にできる。
外側平方関係を$x$と比較し、$T+Tv$の独立性を用いると
\[
 (a-a^2)\ell_x=0,\qquad(a^{-1}-a^2)\ell_y=0
\]
を得る。ここで$T+Tv$に非自明な従属があれば$v\in Q(T)$、
従って$z^2\in Q(T)$となって同じ補題に反する。
両係数は原始6乗根で非零なので$\ell=0$である。
さらに$q_2$を$x,y$と比較して
\[
 q_2x=a^2xq_2,\qquad q_2y=a^{-2}yq_2
\]
を得る。$a^2$は原始3乗根なので、第一式は$q_2$を$ky^2$へ、
第二式は$kx^2$へ制限する。従って$q_2=0$となる。
この場合は先頭環そのものであり、全正常核$ky^3$も因子化する。
\end{proof}

\texttt{verify\_cubic\_positive3\_diagonal\_resonance.py}はskew derivationの四条件と
内部二列、原始6乗根で用いる全非消滅・最大公約数、および$a=-1$標準形での
三中心二次式の全交換子を元の二次関係の次数3イデアル内で厳密検算する。
この検算では補正係数で割らず、任意係数でのAS4存在も仮定しない。
